% ======================================================================
% SECTION 6: HR 9506 - American Physical AI Oncology Leadership Act
% of 2026
% New Federal Authorizing Statute Adapting FDORA Section 3209 and
% 21st Century Cures Act 2016.
% ======================================================================

\hspace{0.5cm} HR 9506 is a \textit{New Federal Authorizing Statute Adapting FDORA
Section 3209 (PL 117-328) and the 21st Century Cures Act of 2016
(PL 114-255)}, drawing the chain of prior authority from the in
silico authorization in FDORA through the interoperability and
information-blocking framework of the Cures Act
\cite{rank4_fdora_fda_modernization_2022, rank1_cures_act_2016}.
The bill is the project brief's "United States Number 1" statute and
absorbs the cross-cutting paper-level material from the prior
template's Implementation Metrics, Discussion, Limitations, and
Conclusions sections. The respectful FDA framing throughout this
section is intentional: the bill is an authorization, not a critique.

The current legislative gap closed by HR 9506 is the absence of a
standalone federal authorizing statute for sustained American
leadership in medical AI and robotics oncology. FDORA Section 3209
already authorizes in silico and computer-model alternatives to
animal testing, mandates diversity action plans, and mandates the
FDA decentralized-trial guidance \cite{rank4_fdora_fda_modernization_2022,
fda_dct_guidance_2024}. HR 9506 extends that authorization from
animal-testing alternatives to physical-AI executor selection and
codifies the American Leadership Index that anchors the cross-bill
implementation timeline.

\subsection{Cross-Bill Implementation Timeline}

\hspace{0.5cm} The American Leadership Index publication cycle is anchored to the
following cross-bill timeline. Use the diagram below as the canonical
schedule for the seven bills (HR 9501 through HR 9507).

\begin{verbatim}
+--------+----------+------------------------------------------------+
| Month  |  Bill    |  Federal Deadline                              |
+--------+----------+------------------------------------------------+
|   12   |  HR 9501 |  Federal API for self-selection                |
|   12   |  HR 9502 |  Per-robot USL transparency reports            |
|   12   |  HR 9506 |  First American Leadership Index publication   | 
|   12   |  HR 9507 |  CMS-billable EHR self-custody endpoint        |
|   18   |  HR 9503 |  Per-protocol modifiable steps matrix          | 
|   18   |  HR 9504 |  Per-task human and AI error rates             |
|   24   |  HR 9501 |  Full state-level self-selection adoption      |
|   24   |  HR 9505 |  Patient-sponsor direct channel coverage       |
|   24   |  HR 9506 |  FDA RTCT pilot expansion (target 25)          |
|   24   |  HR 9507 |  State-level data self-custody harmonization   | 
|   36   |  HR 9506 |  All 50 states with at least one Physical AI   |
|        |          |  oncology trial site authorized                |
|   84   |  HR 9506 |  Seven-year sustained federal investment       |
|        |          |  period 2026 through 2032 ends                 |
+--------+----------+------------------------------------------------+
\end{verbatim}

\subsection{Findings and Declarations}

\hspace{0.5cm} (a) The United States is the global leader in clinical-trial
infrastructure today and must remain Number 1 in the world for
patient benefit during the medical AI and robotics revolution
\cite{kawchak_2026_19244918, main_repo}.

(b) The FDA Real-Time Clinical Trials announcement of April 28,
2026 and the FDA Oncology Center of Excellence Advancing Oncology
Decentralized Trials directory provide evidence of FDA leadership
posture that this bill formalizes into statute
\cite{fda_rtct_press_2026, fda_oce_advancing_oncology_dct_2024}.

(c) FDORA Section 3209 already authorizes in silico alternatives;
HR 9506 extends that authorization from animal-testing alternatives
to physical-AI executor selection and to the federated-learning
infrastructure across CMS-billable sites
\cite{rank4_fdora_fda_modernization_2022}.

(d) The 21st Century Cures Act (PL 114-255) supplies the
interoperability and information-blocking framework that the
American Leadership Index reports against
\cite{rank1_cures_act_2016}.

(e) The author's prior repository deposits document the technical
substrate at site, sponsor, patient-journey, patient-instructions,
USL, and ICH-E6(R3) levels and supply the verifiable basis for the
American Leadership Index reporting metrics
\cite{kawchak_2026_19176370, kawchak_2026_19119939,
kawchak_2026_19396256, kawchak_2026_18810541,
kawchak_2026_18778219, kawchak_2026_18973368}.

(f) Foundation-model evidence (Virchow), trial-matching evidence
(TrialGPT, PRISM, TrialMatchAI), and prediction evidence (SCORPIO,
PROGPATH, the four-simulation paper at
\texttt{new-trial/national-24-7-trial/paper/full-paper/}) anchor the
bill's authorization for investment in foundation models
across pathology, prediction, and trial matching
\cite{vorontsov2024virchow, jin2024trialgpt, gupta2024prism,
abdallah2026trialmatchai, kawchak_2026_19994945}.

\subsection{Definitions}

\hspace{0.5cm} (a) \billterm{Medical AI and robotics revolution} means the period
from approximately 2024 through 2030 during which Physical AI
systems and large-context AI agents (for example Claude Code Opus
4.7 with 1M token context, see \cite{claude_opus_47}) are entering
routine clinical-trial deployment. The four-simulation evidence
base at
\texttt{new-trial/national-24-7-trial/paper/full-paper/sections/results.tex}
provides the operational reference \cite{kawchak_2026_19994945}.

(b) \billterm{American leadership} means a documented United States
lead in:
(i) the number of authorized Physical AI oncology trial sites;
(ii) the number of patient-self-selected trial bookings per year;
(iii) the cumulative documented error-rate reduction across CMS-
billable sites;
(iv) the number of patients whose data is governed under HR 9507
self-custody; and
(v) the number of FDA RTCT pilot trials with patient-sponsor
direct channels under HR 9505.

(c) \billterm{Federal investment} means dedicated appropriations to
the Food and Drug Administration, the National Institutes of
Health (PAR-25-170 and successor programs, see
\cite{nih_par25_170_dht_endpoints}), the Agency for Healthcare
Research and Quality, and the Office of the National Coordinator
for Health Information Technology for the seven-year period 2026
through 2032.

(d) \billterm{American Leadership Index} means the public report
published annually by the FDA Oncology Center of Excellence with
the five metrics defined in 6.3(b) above.

\subsection{American Leadership Authorizations}

\hspace{0.5cm} (a) Authorization for the FDA RTCT pilot program to expand from the
initial two trials (TRAVERSE, STREAM-SCLC) to at least 50 trials by
2028 \cite{fda_rtct_press_2026}.

(b) Authorization for the FDA Oncology Center of Excellence to
publish per-state dashboards of Physical AI oncology clinical trial site
authorization counts, harmonized with the prior trial site documentation
framework at
\texttt{new-trial/site/05-state-regulations/} and at
\texttt{national-platform/new\_trial\_psl/05\_title22\_regulations.tex}
\cite{kawchak_2026_19176370, kawchak_2026_19244918}.

(c) Authorization for federated-learning infrastructure across CMS-billable sites, drawing on the framework at
\texttt{national-platform/federated\_learning/} (chunks 1 through 4)
\cite{kawchak_2026_19244918}.

(d) Authorization for federated investment in foundation models for
pathology (Virchow), prediction (SCORPIO, PROGPATH), and trial
matching (TrialGPT, PRISM, TrialMatchAI), with annual per-model
reporting on patient-benefit metrics \cite{vorontsov2024virchow,
jin2024trialgpt, gupta2024prism, abdallah2026trialmatchai}.

(e) Authorization for a public American Leadership Index published
annually with the five metrics defined in 6.3(b).

\subsection{New Physical-AI Reversal Guardrail}

\hspace{0.5cm} A new physical-AI guardrail introduced by this paper is anchored
under HR 9506's annual American Leadership Index reporting
authority: physical-AI executor selections and procedural
modifications must always permit a patient-initiated reversal to
human-only execution within the same calendar day. The operational
baseline is \texttt{patient-journey/master\_journey.py} and the manager at
\texttt{sponsor/final\_paper/scripts/core\_agents/audit\_trail\_manager.py}
\cite{kawchak_2026_19119939, kawchak_2026_19396256}.

\subsection{Implementation}

\hspace{0.5cm} (a) The Secretary of Health and Human Services, acting through the
Food and Drug Administration, the National Institutes of Health,
the Agency for Healthcare Research and Quality, and the Office of
the National Coordinator for Health Information Technology, shall
publish the first American Leadership Index within 12 months of
enactment.

(b) The FDA RTCT pilot shall expand to at least 25 trials within 24
months of enactment, with the target of 50 trials by 2028
\cite{fda_rtct_press_2026}.

(c) 
\texttt{national-platform/new\_template/sections/implementation\_strategy.tex} implementation pattern
informs the implementation guidance \cite{kawchak_2026_19244918}.

\subsection{Reporting and Enforcement}

\hspace{0.5cm} (a) Annual American Leadership Index report.

(b) Congressional review under FDORA authorities; explicit
alignment with the FDORA Section 3209 in silico authorization
\cite{rank4_fdora_fda_modernization_2022}.

\subsection{Future Work Tracks}

\hspace{0.5cm} HR 9506 authorizes federal investment that supports three future
tracks for delivering all seven bills' obligations.

Track A (Single big model performing all bills' obligations). A
future generation of Claude Code Opus 4.7 (or equivalent
foundation-model agent) holds all seven bill texts plus the
AVAILABLE DIRECTORIES content in a single 1M-token context and
emits implementation deliverables (machine-readable eligibility
reports, USL transparency reports, modifiable-steps matrices,
error-rate publications, real-time aggregator endpoints, the
American Leadership Index, and self-custody endpoints) in a single
inference loop. This track is the natural extension of the four-simulation paper at
\texttt{new-trial/national-24-7-trial/paper/full-paper/}
\cite{kawchak_2026_19994945, claude_opus_47}.

Track B (Distributed agents per bill, coordinated by a smaller
orchestrator). Each of the seven bills has its own dedicated agent
(self-selection agent for HR 9501, robot-choice agent for HR 9502,
modification agent for HR 9503, error-rate agent for HR 9504,
sponsor-channel agent for HR 9505, leadership-index agent for HR
9506, self-custody agent for HR 9507), all coordinated by a
smaller orchestrator that runs on hardware as constrained as the
Core i5-6200U with 4 GB RAM at
\texttt{sponsor/final\_paper/168\_hours/instructions/core\_i5\_6200u\_4gb/README.md}
\cite{kawchak_2026_19396256}.

Track C (Patient-side native deployment). A future generation
deploys a single integrated patient-side agent that operates HR
9501 through HR 9507 from the patient's own device, using
patient-controlled credentials and HIPAA-compliant authentication
to interact with provider, payer, and sponsor endpoints. This
track is the natural extension of
\texttt{patients/patient\_robot\_instructions\_fixed.tex} pages 1
through 10 plus the digital-twin infrastructure at
\texttt{patient-journey/stage\_03\_digital\_twin.py}
\cite{kawchak_2026_18810541, kawchak_2026_19119939}.

\subsection{Future Work Deliverables}

\hspace{0.5cm} HR 9506 authorizes federal investment supporting three concrete
future deliverables:
(i) a TRIPOD+AI / CREMLS-compliant retrospective validation of HR
9504's error-rate baseline on a real-patient cohort, drawing on the
SHIELD-RT design pattern;
(ii) a public benchmark of HR 9501's machine-readable eligibility
report performance across the TrialGPT, PRISM, and TrialMatchAI
matching systems on a shared dataset; and
(iii) an FDA-aligned RTCT pilot submission that uses Simulation 4
(\texttt{sponsor/final\_paper/168\_hours/}) as the sponsor-side input
and HR 9505's patient-sponsor direct channel as the patient-side
input \cite{fda_rtct_press_2026, kawchak_2026_19994945,
kawchak_2026_19396256, kawchak_2026_19119939}.

\subsection{Persistent Themes Across the Seven Bills}

\hspace{0.5cm} The persistent themes anchored under HR 9506's American Leadership
Index are: (1) Adaption / Revision framing, with every bill
explicitly stating the prior United States legislative document it
builds from; (2) respectful FDA framing, with every bill positioning
the FDA and other governing bodies as enabling rather than gating
while remaining opportunistic under the patient's control based
primarily on the patient's own interests; (3) individual patient
plus broad-adoption framing, with every bill naming per-patient
examples (PAT-2026-0042 from \texttt{patient-journey/}) AND broad-
adoption deadlines (12, 18, 24 months across CMS-billable sites and
50 states); and (4) United States Number 1, with every bill closing
with the leadership framing required to keep the United States
Number 1 in the world for patient benefit during the medical AI and
robotics revolution.

\subsection{Implications for Cancer Patient Control}

\hspace{0.5cm} Large language model-scale repository context with one million tokens capacity demonstrated in
\texttt{new-trial/national-24-7-trial/paper/full-paper/}, the four
LLM simulations summarized at \cite{kawchak_2026_19994945}, the
168-hour sponsor extension at \texttt{sponsor/final\_paper/168\_hours/},
the 10-stage single-patient journey at \texttt{patient-journey/},
and the 10-robot-category patient instructions at
\texttt{patients/patient\_robot\_instructions\_fixed.tex} together
show that the technical substrate is ready. The seven bills supply
the missing legal substrate. Together they enable the cancer
patient to take control of the patient's disease in ways that
prior error-prone and physically and intellectually limited human-only care could not support.

\subsection{Discussion and Limitation}

\hspace{0.5cm} The bill is the most ambitious of the seven; it assumes seven-year
sustained federal investment which is subject to appropriations
cycles. The American Leadership Index relies on year-over-year
metric stability across 50 states. Across the seven bills, the
draft language is by an independent author and has not been
reviewed by Congress, the FDA, or any State legislature. The bills
are not endorsed by Anthropic, OpenAI, Google, or any technology
vendor. The patient-control framework remains an operational
hypothesis; outcome studies are required to confirm that the
bills, taken together, increase participation and decrease error
rates without unintended consequences. The bills do not address
pediatric or vulnerable-population protections beyond the existing
21 CFR 50 Subpart D at
\texttt{national-platform/21cfr50\_adapt/02\_irb\_review\_pediatric.tex}
\cite{rank4_fdora_fda_modernization_2022, kawchak_2026_19040707}.

\subsection{Forward Path}

\hspace{0.5cm} The next Claude Code Opus 4.7 1M Max generation pass will refine
this paper with additional ASCII diagrams, additional individual
patient and robot examples drawn from the AVAILABLE DIRECTORIES, and
additional cross-bill cross-references; the author will then
perform additional senior-author white-space and table-column-width
formatting passes to remove any orphans, widows, or large empty
page regions before publication. References anchor the forward
path: \cite{kawchak_2026_20045457}, \cite{fda_rtct_press_2026},
\cite{kawchak_2026_19994945}, \cite{kawchak_2026_19396256},
\cite{kawchak_2026_19244918}, and \cite{kawchak_2026_18810541}.

\subsection{Patient-Control Framings}

\hspace{0.5cm} For the individual patient, a cancer patient anywhere in the
United States can verify, in the published American Leadership
Index, that the patient's State has at least one authorized
Physical AI oncology trial site, that the per-task error rates have
improved year-over-year, and that the FDA RTCT pilot covers the
patient's cancer type. PAT-2026-0042 in \texttt{patient-journey/}
illustrates the per-patient outcome trajectory across the 10-stage
journey \cite{kawchak_2026_19119939}.

For broad adoption, all 50 states must have at least one authorized
Physical AI oncology trial site within 36 months of enactment; the
United States retains Number 1 status in the world for patient
benefit during the medical AI and robotics revolution by virtue of
the per-state, per-cancer, per-site dashboard publication; federal
appropriations sustain the FDA, NIH, AHRQ, and ONC infrastructure
for seven years \cite{kawchak_2026_19244918, kawchak_2026_19994945,
kawchak_2026_19396256}.
